\documentclass[11pt]{article}
\usepackage{makeidx}
\usepackage{amssymb}
\usepackage{amsfonts}
\usepackage{amsmath}
\usepackage{setspace}
\usepackage[hmargin={1.2in,1.2in},vmargin={1.25in,1.25in}]{geometry}
\usepackage{graphicx}
\usepackage[shortlabels]{enumitem}
\usepackage{mathpazo}

\setlength{\parskip}{.25cm}
\setcounter{secnumdepth}{3}
\graphicspath{{./graphics/}}

\begin{document}

\thispagestyle{empty}
\setstretch{1}
\setlength{\parindent}{0pt}


\begin{flushright}
	\today
\end{flushright}

\bigskip

Dear Yuriy,

Thank you for giving me the opportunity to review this paper for the \textit{American Economic Review}, and apologies for the slight delay in submitting my report. I have now read the paper and have two main comments:

\textbf{1. Identification}: The empirical analysis relies on the identification assumption that the growth in the service flow of housing does not differ between extending and non-extending properties. While the authors provide solid evidence to support this assumption, there are areas that could benefit from further clarification. Specifically, I am concerned about potential geographical biases, the impact of high-income tenants on property values, and the effect of various taxes on lease extension decisions.

\textbf{2. Interpretation of \( y^* \)}: The authors decompose the expected long-term housing yield into three components: the long-term safe asset yield (\( r^* \)), long-term risk premia (\( \zeta^* \)), and long-term dividend growth (\( g^* \)). They claim that the decline in \( y^* \) is primarily due to a decline in \( r^* \), consistent with broader trends in long-term interest rates. However, comparing \( y^* \) to 10-15 year forward rates (as the authors do in the paper) may not be appropriate, as \( y^* \) is expected to prevail 70-160 years into the future. Considering ultra-long-run bonds, like Austria's century bond, for example, might provide more robust insights into post-pandemic trends. Alternatively, the authors could report an estimate of \( y^* \) based on properties with a duration that is comparable to more policy-relevant horizons (e.g., 10 to 20 years).

You can find a more detailed explanation of these points in my report, along with a few additional comments.

In sum, the paper presents some solid empirical evidence on an interesting object, the long-run housing yield. However, I have some doubts about its policy relevance. For these reasons, I recommend a weak \textit{Revise and Resubmit}. I believe the paper could have significant potential contingent on being able to speak more directly to \( r^* \).

I hope this letter, along with my report, will help you in making a decision on the manuscript.

Best regards,\newline
Ambrogio Cesa-Bianchi

\end{document}

Also, I would have benefited from a theoretical framework to understand the lease extensions and the link between $y^*$ and $r^*$. But there is obviously already a lot in the paper.
